\section{What Is Logic?}

Logic is a method for systematically evaluating arguments of various
types. Although this course is a study in \textit{formal} logic, the
main concepts are best introduced informally.

\subsection{Arguments} \cite[][Chapter 1, pp. 2-6]{Mag25}

An \textit{argument} is a sequence of sentences. It comprises some
\textit{premises} and a \textit{conclusion}.

\begin{argument}
\item[] My shoes are either in the den or in the bedroom.
\item[] My shoes are not in the den.
\item[\therefore] My shoes are in the bedroom.
\end{argument}

In real-world examples, the conclusion can occur anywhere in such sequence.

\begin{argument}
\item[] My shoes are either in the den or in the bedroom. They must be in
  the bedroom, because they are not in the den.
\end{argument}

But when we turn to formal logic, conclusions will be found at the end.

Informal arguments will often use special words to indicate which sentences
are meant to be premises and which is the conclusion.

\begin{argument}
\item[] \textbf{Conclusion indicator words:}
  \textit{so, therefore, hence, thus, accordingly, consequently, ... }
\item[]  \textbf{Premise indicator words:}
  \textit{since, because, given that, ... }
\end{argument}

The premises and conclusions of these informal arguments are \textit{sentences},
but not just any type of sentence will do. We are not interested in commands or
questions, but only \textit{declarative} sentences. These are the sentences that
we can assign a \textit{truth value}; sentences that we can call \textit{true}
or \textit{false}.

\subsection{The scope of logic} \cite[][Chapter 2, pp. 7-17]{Mag25}

In arguments like the one outlined above, we note that the conclusion is
connected to the premises in a fundamental way. The conclusion is a
\textit{consequence} of the premises; whenever the premises are true, the
conclusion is also true. When the consequence of an argument follows in this way
from its premises, we will say that the argument is \textit{valid}. Otherwise, it
is \textit{invalid}.

Although we will keep the discussion informal for the time being, note that we can
talk about arguments by representing their premises and conclusions using capital
letters (with or without subscripts). Using this convention, we can quickly define
this notion of consequence (\textit{entailment}):

\begin{argument}
\item[] The sentences $A_1, \dots, A_n$ \textit{entail} the sentence $B$ if and
  only if $B$ is never true when all of $A_1, \dots, A_n$ are true.
\end{argument}

We need to flesh out 'never' a bit more in our formal defnition, but this gives us
a sense of what we mean by an valid argument and its components. From now on we
will use 'iff' to mean 'if and only if'.

\begin{argument}
\item[] An argument is \textit{valid} iff its premises entail its conclusion.
\item[] An argument is \textit{invalid} iff it is not valid.
\end{argument}

An invalid argument can be detected by noting a \textit{counter-example}: an
example of an argument with the same form where the premises are all true but the
conclusion is false.

Arguments can be valid in a number of different ways. These depend on the limits we
place on the types of counter-examples that we will accept and invalidating.

\begin{argument}
\item[] Nomologically valid: must respect the laws of nature
\item[] Conceptually valid: must respect conceptual relations between words
\end{argument}

The replacement of sentences by letters in valid arguments is our first step to
formal logic. We understand \textit{formal} as pertaining to the \textit{form} of
such arguments, not as a synonym for \textit{fancy}. One consequence is that many
conceptually valid arguments cannot be formally validated. So this formal argument
is valid:

\begin{argument}
\item[] Either X or Y.
\item[] Not Y.
\item[\therefore] X.
\end{argument}

while this one is not, unless we make clear the premise that $X iff Z$:

\begin{argument}
\item[] Either X or Y.
\item[] Not Y.
\item[\therefore] Z.
\end{argument}

Note that the dependence on the forms of arguments is fundamental here. We can have
both valid arguments with untrue conclusions and invalid arguments whose premises
and conclusions are all true.

\begin{quote}
  The important thing to remember is that validity is not about the \textit{actual}
  truth or falsity of the sentences in the argument. It is about whether it is
  \textit{possible} for all the premises to be true and the conclusion to be not
  true at the same time. \cite[][p15]{Mag25}
\end{quote}

Logic is exceptional at detecting invalid arguments but cannot help at all with
assessing the truth of premises. Additionally, formal logic is limited to the
assessment of \textit{deductive} arguments. It does not help with the assessment
of the many inductive arguments that populate empirical sciences and everyday
experience.

\begin{argument}
\item[] The sun has risen every day so far.
\item[\therefore] The sun will rise tomorrow.
\end{argument}

While this is a perfectly reasonable inference to make, it is not a deductivley
valid argument.

\subsection{Other logical notions} \cite[][Chapter 3, pp. 18-25]{Mag25}

We can also define some other properties of sentences and relations between them:

\begin{argument}
\item[] Sentences that are always true are \textit{necessary truths}.
\item[] Sentences that are never true are \textit{necessary falsehoods}.
\item[] Sentences that are neither necessary truths nor necessary falshoods are
  \textit{contingent sentences}; they are sometimes true and sometimes false.
\item[] Pairs of sentences that have the same truth value in every situation are \textit{necessarily equivalent}.
\item[] Collections of sentences that are all true together in some situation are \textit{jointly possible};
  those that are not are \textit{jointly impossible}.
\end{argument}

